\chapter{Introduction}
\label{ch:introduction}

\section{Motivation}

Ant scans are valuable because they can turn specimens that are difficult to
measure repeatedly into a source of comparable three-dimensional traits. Once
anatomically corresponding models are available, quantities such as head
width, Weber's length, and appendage proportions can be extracted
systematically rather than measured by hand under a microscope. The practical
promise is therefore not only better mesh fitting, but a way to study
biodiversity at population scale.

In the first week of July 2026, ten preprocessed ant scans were passed through the
registration pipeline and the resulting loss curves were inspected. All ten
descended cleanly and converged without incident; by the only quantitative
signal the system produced, every run had succeeded.

Visual inspection told a different story. Four of the ten fitted meshes were
not recognisable ants. The template had been deformed into a smooth mass that
approximately matched the target's silhouette and enclosed volume, but bore no
anatomical correspondence to it: legs had collapsed into the gaster, and the
head had been absorbed into the thorax. Critically, the loss on several of
these four runs was \emph{lower} than on the runs that had visibly succeeded.

\findingmark{Established here; recurs throughout this report as \F{F1}.}%
This discrepancy is the starting point for the work reported here. The
objective minimised during fitting is a symmetric point-to-point distance
between two surfaces. It quantifies whether the fitted mesh \emph{covers} the
target, not whether vertex $i$ of the template has landed on the anatomical
point that vertex $i$ is defined to represent. These properties can diverge;
when they do, the failure produces no exception, no divergent loss curve, and
no indication in the log that anything is wrong.\footnote{This is recorded as
\F{F1} in the findings ledger (Section~\ref{sec:ledger}) and motivates the
correspondence measurement developed later in the investigation.}

This contradiction is the central problem addressed by the internship. In a
registration method based on numerical optimisation, the objective function is
not meant to be a visual proxy for biological correctness; it is only one
surrogate of the desired outcome. When that surrogate rewards global coverage
without enforcing anatomical identity, the algorithm can converge to a strong
but biologically meaningless solution. The issue is therefore not merely one of
optimisation failure but of specification failure: the system was measuring the
wrong thing, or at least an incomplete version of it.

The project began from a fundamental question: how can a parametric
mesh be fitted to imperfect animal scans if the pipeline has no objective that
disambiguates good correspondence from a plausible but wrong deformation? The
internship addressed this question by moving from a mesh-quality-centred
hypothesis to a more precise investigation of initialisation, correspondence,
and constrained optimisation. A method is not considered reliable simply
because it reaches a numerically stable minimum, but because it reaches a
biologically meaningful one.

\section{Scope of the internship}

The scope of the internship was intentionally defined around a concrete
technical bottleneck rather than a broad, open-ended exploration. The initial
assignment list pointed to preprocessing, segmentation, neural initialisation,
and hybrid registration as the main design axes. In practice, the early
experiments showed that the dominant bottleneck was not the raw mesh quality
alone, but the interaction between poor initialisation and a weak sense of
anatomical correspondence. This insight re-centred the work away from a
preprocessing-only narrative and toward the underlying determinants of
registration quality inside the optimiser itself.

The resulting investigation therefore combined three layers of work. First, it
examined the assumptions of the project and tested them against real data.
Second, it developed measurement techniques capable of distinguishing surface
coverage from true anatomical correspondence. Third, it evaluated candidate
interventions against explicit criteria and retained only those that improved
the pipeline in a demonstrable and reproducible way. This report follows that
sequence because the evidence gathered at each stage changed the direction of
research and justified the choices that were ultimately retained in the final
SMILify pipeline.

The internship brief, reproduced in the technical data sheet appended to this
report, lists six assignments: unify the preprocessing pipeline, stabilise a
Shape-Diameter-Function-based segmentation, integrate neural initialisation,
design a hybrid registration framework, benchmark the result, and deliver
documentation. Read in sequence, this list implies a specific hypothesis about
the source of registration failure: that the input meshes are the limiting
factor, and that cleaning them will improve the fits. Section~\ref{sec:turn}
describes how an early, direct test of that hypothesis produced a null result,
and how the project was redirected on the basis of that evidence. The present
chapter sets out why animal shape modelling is a materially harder problem than
its human counterpart, and what state the project was in at the outset.

\section{Parametric shape modelling for non-human subjects}

Parametric mesh registration deforms a rigged, segmented template mesh until it
matches a target scan, governed by a shape coefficient vector $\bm\beta$, joint
rotations $\bm\theta$, and a global translation $\bm\gamma$. A correctly fitted
template does more than resemble the target visually: because every vertex
retains a fixed anatomical identity across all fits, the result places
otherwise unrelated scans into a shared coordinate system in which
measurement, comparison, and population-level analysis become possible. This
property, rather than visual resemblance, is what makes the technique useful.

For humans, the underlying recipe is well established. \term{SMPL} (Skinned
Multi-Person Linear)~\cite{smpl} was learned from thousands of registered scans
of cooperative human subjects. Animal
subjects violate the assumptions this recipe depends on: shape variation across
species substantially exceeds variation within a human population, the
subjects cannot be instructed to hold a pose, and a representative sample
cannot practically be brought into a scanning facility. Zuffi et
al.~\cite{smal} addressed this by learning a shape space, \term{SMAL} (Skinned
Multi-Animal Linear), from
scans of toy figurines rather than live animals, and demonstrated that the
resulting model generalises to photographs of real animals, including species
absent from the training set.

\term{SMILify} extends this approach to insects, and in doing so removes a
further set of assumptions. Where a quadruped presents four limbs, an ant
presents six, plus two antennae and a pair of mandibles, all thin, all
articulated, and several represented by only a few dozen vertices in the
template mesh. In ant morphology, the thorax is more precisely the mesosoma;
the former term is retained here when it refers to the model's part naming.
The scan data originate from \term{scAnt}, an open-source macro
photogrammetry scanner~\cite{scant}, and from micro-computed tomography.
Meshes produced this way commonly contain holes, self-intersections, arbitrary
global orientation, non-manifold geometry, and missing surface exactly at the
structures of greatest anatomical interest. The technical data sheet for this
internship states the resulting problem directly: existing registration
pipelines are designed for clean, watertight meshes in near-canonical poses,
and fail or become unstable on input of this kind.

\section{Project state at the outset}

\term{SMILify} is built directly on \term{SMALify}~\cite{smalify}, Biggs et
al.'s optimisation-based framework for fitting the SMAL quadruped model to
video and image data, which is itself built on the original SMAL
model~\cite{smal}. The repository's own framing of its ambition is direct: to
turn any rigged mesh into a SMAL-compatible parametric model, regardless of
skeletal topology, with \term{SMIL} (Skinned Multi-Insect Linear) as the
current insect-focused instance of a more general goal.
\footnote{\url{https://github.com/FabianPlum/SMILify}, README, accessed
September 2026.}

The repository on 1 July 2026 was an active codebase rather than a blank
starting point. \code{fitter\_3d/} contained the mesh registration pipeline
used to author new parametric models by fitting a template to a collection of
target \code{.obj} scans; \code{custom\_processing/} held model-preprocessing
helpers, including mesh registration utilities and beta regressors;
\code{3D\_model\_prep/} held the Blender addon used to build and edit
parametric \code{.pkl} models; \code{config.py} held the legacy root
configuration still required by \code{fitter\_3d}. Fitting jobs were, and
remain, invoked directly against a YAML configuration, for example:
\begin{quote}\small\ttfamily
python -m fitter\_3d.optimise --mesh\_dir <target\_meshes>\\
\hspace*{1em}--yaml\_src fitter\_3d/ants\_cfg.yaml
\end{quote}
A Shape Diameter Function module intended for part-aware segmentation had not
been validated on incomplete meshes. \term{PointNet}-based predictors for pose
and shape initialisation were implemented but unevaluated. The available scan
corpus itself indicated the scale of the data-quality problem: of roughly 380
raw scans, approximately 250 survived processing and 86 were curated as clean
enough for registration, meaning that more than three-quarters (approximately
77\%) of the collected material was discarded before registration was ever
attempted on it.

The repository's own issue tracker records a longer-term ambition beyond this
internship's scope: building a SMIL model from the full antscan collection
rather than a curated subset, tracked as a future objective referred to
internally as \term{SMILify} Gen2. The registration accuracy work reported
here, and in particular the elimination of mesh-quality as the dominant
bottleneck (Section~\ref{sec:turn}), is a direct prerequisite for that objective,
since a corpus-scale shape space can only be trusted if the registrations
feeding it are individually reliable.

\section{The redirection of the project}
\label{sec:turn}

\findingmark{The result that redirected the project; recorded as \F{F2}.}%
Consistent with the hypothesis implicit in the brief, the first weeks of the
internship were spent building a unified preprocessing pipeline, with explicit
diagnostics, reproducibility, and quality classification. That pipeline is
retained in the delivered system because these outputs make later fitting
decisions auditable. It was then evaluated directly against the existing
preprocessing pipeline in a controlled comparison, holding the registration
stage fixed. The result was a statistical tie: mean Chamfer distance of
$0.000771$ for the new pipeline against $0.000737$ for the existing one
(lower is better). The nominal difference favoured the existing pipeline but
did not reach significance (\needcheck{insert $n$, paired-test statistic, and
$p$-value}). Cleaner input meshes did not produce more accurate
registrations.

The null result was informative because it tested the location of the
bottleneck rather than merely the performance of one implementation. A
positive result would have supported further preprocessing; this result ruled
out mesh quality as the dominant explanation and redirected the internship
toward the question of what was governing registration accuracy.

Answering that question required an instrument that did not yet exist: a
measurement of anatomical correspondence, as distinct from surface coverage.
This was constructed from synthetic data generated by the model itself:
parameters are sampled from the model's own distributions, the template is
posed accordingly, the result is exported as an anonymous mesh, optionally
degraded, and passed back through the fitter. Because the mesh originates from
the model, the true identity of every vertex is known by construction, and
per-vertex correspondence error becomes directly measurable rather than
inferred indirectly from surface distance.

This instrument produced two results that determined the shape of the
remaining investigation. First, it localised the correspondence error:
approximately 83\% of it does not cross a part boundary (\F{F3}), indicating
that points are more often placed incorrectly \emph{within} the correct
anatomical part than assigned to an incorrect part altogether. This finding
rules out an entire family of proposed interventions aimed at improving
part-level discrimination, since such interventions target the remaining
17\%. Second, fitting from the ground-truth pose, rather than from a default
initialisation, recovers most of the lost accuracy, indicating that the
dominant failure mode is the optimiser becoming trapped in a poor local
minimum near a bad starting point, rather than the objective itself being
poorly shaped (\F{F4}). These two findings jointly reoriented the remaining
work toward initialisation and within-part correspondence, and away from
interventions aimed at coarse mesh quality or part-level segmentation.

The same audit also exposed a structural limitation in the sampling scheme:
thin distal structures receive far less vertex support than proximal body
parts, independently of pose (\F{F5}). In the evaluated template this meant
38 vertices for the tarsus and 8 for the pretarsus, compared with hundreds on
the proximal body. This observation does not by itself identify a successful
remedy, but it explains why a global surface metric can underweight precisely
the structures that matter most anatomically.

The remainder of the internship proceeded by systematic elimination.
Candidate interventions, including robust losses, correspondence changes, and
sampling strategies, were each tested against a criterion registered in
advance of the experiment; the complete elimination map is preserved in the
project diagnostics. The interventions retained for SMILify are joint rotation
limits consumed as a fitting prior, per-joint scale barriers validated on a
separate 838-specimen corpus (757 worker specimens plus 81 clean specimens),
and a staged hierarchical fitting recipe that outperforms the original
single-pass fitter on every integrity metric measured. The 838 specimens are
not the 86 curated scans described above; they belong to the larger validation
run.

The closing phase of the internship shifted from registration accuracy to its
downstream application. A correctly registered mesh is not an end in itself
but an instrument: once every specimen occupies a shared parameterisation,
measurements that entomologists currently extract manually under a
microscope, including head width, Weber's length, and related proportional
relationships, become extractable automatically and at population scale.
The validation work evaluates whether the delivered fits are accurate enough
to support this claim, including an explicit treatment of where the validation
sample is not representative and where a reported correlation is measured on
data selected by the system under evaluation.

\section{Established results}
\label{sec:ledger}

The results below recur throughout this report and are registered once here;
each is cited by identifier (\F{F1}, \F{F2}, and so on) at the point where it
is used. Full evidence for each is preserved in the associated diagnostics
and experiment records.

\printfindingledger

\section{Contributions}

The results of this internship are summarised below, grouped by their level of
confidence.

\paragraph{Delivered and validated.} A correspondence-accuracy measurement
built on synthetic ground truth, where no equivalent measurement previously
existed; joint rotation limits consumed as a fitting prior in the 3D
registration path, closing the one fitting path in \term{SMILify} that had
previously optimised pose without constraint; per-joint scale barriers,
validated on a separate corpus of 838 specimens and reducing anterior scale
pathology (unrealistic scaling of anterior body groups) by 58--70\% on the
affected outlier tail; a unified, configurable
preprocessing module with diagnostic output and quality classification; a
staged hierarchical fitting recipe that outperforms the original single-pass
fitter on every integrity metric measured.

\paragraph{Established.} The five findings listed in Section~\ref{sec:ledger}, each
independently reproduced, which collectively relocate the primary bottleneck
of the registration pipeline from mesh input quality, as assumed in the
original brief, to initialisation and within-part correspondence.

\paragraph{Eliminated.} A documented set of candidate interventions tested
against pre-registered criteria and rejected, each with the specific evidence
that determined the outcome. The associated diagnostics remain available so
that subsequent work is not repeated unnecessarily.

\paragraph{Open.} Three problems that this investigation localised but did
not resolve: initialisation without prior knowledge of the target pose,
within-part correspondence on thin articulated geometry, and systematic
under-representation of distal structures during sampling. These are stated
explicitly as open problems rather than presented as partial successes.

\section{Structure of this report}

The report begins with the motivation and technical background for parametric
registration of insect scans. It then situates the project within the host
organisation, the computational environment, and the research context in
which the SMILify work was carried out. The appendix maps the original
internship assignments to the delivered repository. Within this compact
structure, the central evidence chain remains visible: biological need,
observed failure, tested hypothesis, measured bottleneck, and interventions
retained for their value to SMILify.
